\section*{Open Questions}
\addcontentsline{toc}{section}{Open Questions}

Five truly-open questions surfaced by re-reading the paper and by the
replication itself. Each is a candidate research task, not future-work
fluff.

\paragraph{Q1. Empirical value and stability of $\delta$ in Eq.(1).}
\emph{Basis.} Brahme asserts $\delta \approx 0.2$ but never fits it to a
deposited cohort in this paper (C11 asserted, not derived). Our
monotonicity check confirms only that $P^{+}$ rises with $\delta$; the
value drives $P^{+}_{\max}$ by $\sim$5\,pp across $\delta\in[0,1]$ in
our smoke.
\emph{Next steps.} (a) Assemble a public paired
tumor-control / severe-normal-tissue-complication cohort (QUANTEC,
MSK/NIH, EORTC/RTOG released tables). (b) Joint-fit $\delta$ alongside
$(D_{50},\gamma)$ per (site, modality). (c) Compare global $\delta$
vs.\ per-site vs.\ site$\times$modality by AIC/BIC. (d) Test whether
$\delta$ varies with fractionation. (e) Publish fit + posterior.

\paragraph{Q2. Does RHR beat LQ and RCR on external cell-survival data?}
\emph{Basis.} Claim C9 (RHR $>$ LQ for LDHS) is central but its
parameters (a, b, repair rates, NHEJ/HR partitioning) are not deposited
here; live in refs [1--3, 23, 34, 45]. Untested by our audit.
\emph{Next steps.} (a) Pull open cell-survival curves (PIDE, RaySearch,
GDC) covering LDHS-prone and non-LDHS lines. (b) Fit LQ, LQ-L, RCR,
RHR under one optimizer; report AIC/BIC + cross-validated log-likelihood;
predict on held-out $<0.5$\,Gy points. (c) Publish a portable Python
RHR implementation.

\paragraph{Q3. Do reported per-modality $\gamma_C$ values follow from microdosimetric first principles?}
\emph{Basis.} C5 ($\gamma_C\approx 4$ for n/C vs.\ 5--6 for photons/Li,
Fig.\,15) is the paper's only quantitative modality-level claim;
figure-embedded and not recovered by pdftotext.
\emph{Next steps.} (a) Re-parse the PDF with GPU Marker/Nougat to
recover the Fig.\,15 insert. (b) Compute $\sigma_D/\bar{D}$ per modality
from PIDE/PSTAR/SRIM stopping-power tables and ICRU-style cell-size
assumptions. (c) Propagate through
$\gamma_C\approx\gamma_{C,0}/\sqrt{1+(\sigma_D/\bar{D})^2}$
and compare against Brahme's values; report any $>20\%$ mismatch.

\paragraph{Q4. Reproducing BIOART/IVPA on an open PET-CT dataset.}
\emph{Basis.} Claim C12 is the paper's most concrete clinical proposal
(per-voxel $D_{0,\text{eff}}$ from paired pre/mid-treatment FDG-PET-CT)
but is illustrated on a single unnamed lung case with no released
images or maps.
\emph{Next steps.} (a) Identify open datasets with mid-treatment PET-CT
re-imaging at $\approx 18$\,Gy (TCIA HN-PET, NSCLC-Radiogenomics,
public EORTC releases). (b) Implement the uptake-to-$D_{0,\text{eff}}$
mapping per Brahme's derivation; deposit under a permissive license.
(c) Simulate a $D_{0,\text{eff}}$-driven boost and compare against
standard planning via TCP/NTCP surrogates. (d) Concordance study on
follow-up outcome if any dataset supports it.

\paragraph{Q5. Pipeline: how should LUCID handle predatory-venue reviews with no primary artifact?}
\emph{Basis.} This slot exposed a pipeline failure mode --- the paper
passed initial screening as \texttt{candidate\_curated} despite being
a review, in a Gavin Publishers venue, with zero primary artifact. It
consumed an audit slot and required a manual retiering on 2026-06-25.
\emph{Next steps.} (a) Pre-audit venue-quality check
(Cabells/DOAJ-removal/Kscien lists). (b) Static PDF-content check:
count Methods pages, presence of a Data Availability statement, any
parameter table with $>5$ rows. (c) Log any post-audit verdict flips to
tune the gate. (d) Route legitimate reviews to a distinct
\texttt{REVIEW\_AUDIT} track that captures citation-graph and
reachable-primary-refs coverage rather than forcing a REPLICATED
verdict. (e) Publish the gate rules as a standalone package.
